
ce qui
montre que pour \(1\le k\le d\), \(\lambda^k N\) s'envoie sur la \(k\)-ième fonction symétrique des \(N_i\), et par suite que \(f\) est injectif et identifie \(A\) à l'anneau des polynômes symétriques en les \(N_i\).

On a alors \(\lambda_{-1}(N)=\prod_{i=1}^{d}(1-N_i)\), et \(\lambda_t(\lambda_{-1}(N))=\lambda_t\bigl(\prod_{i=1}^{d}(1-N_i)\bigr)\). Les \(\lambda^p(\lambda_{-1}(N))\) sont des polynômes par rapport aux \(N_i\)~; si pour certain indice \(i\) on fait \(N_i=1\), on voit que \(\lambda_t(\lambda_{-1}(N))\) devient \(\lambda_t(0)=1\), donc que les \(\lambda^p(\lambda_{-1}(N))\) s'annulent pour \(N_i=1\), donc qu'ils contiennent \(1-N_i\) en facteur. Donc ils contiennent le produit des \(1-N_i\) en facteur, i.e. ils sont divisibles par \(\lambda_{-1}(N)\) dans \(B\), donc aussi dans \(A\), le quotient de deux polynômes symétriques étant évidemment symétrique.

5.3. Comme \(A\) est intègre, on définit sans ambiguïté un élément \(\lambda^p(N,x)\) en posant (\(p\ge1\)) :
\begin{equation}\tag{5.3.1}
  \lambda^p(N,x)\mathbin{.}\lambda_{-1}(N)=\lambda^p(x\mathbin{.}\lambda_{-1}(N)).
\end{equation}

L'élément \(\lambda^p(N,x)\) est donc un polynôme par rapport aux \(\lambda^i(x)\), \(\lambda^j(N)\), isobare de poids \(p\) par rapport aux \(\lambda^i(x)\), car il en est de même de \(\lambda^p(x\mathbin{.}\lambda_{-1}(N))\) d'après 2.3.

Soient \(d'\) un entier \(\ge d\), et \(A'\) la \(\mathbb Z\)-\(\lambda\)-algèbre engendrée par les générateurs \(N'\) et \(x'\) soumis aux relations \(\lambda^i(N')=0\) pour \(i>d'\)~; d'après la proposition 4.7, il existe un \(\lambda\)-homomorphisme et un seul \(f:A'\longrightarrow A\), envoyant \(N'\) sur \(N\) et \(x'\) sur \(x\). De la définition (5.3.1) résulte que \(f(\lambda^p(N',x'))=\lambda^p(N,x)\) pour tout \(p\ge1\).

Si \(B\) est un \(\lambda\)-anneau, et \(N''\) un élément de \(B\) tel que \(\lambda^k(N'')=0\) pour \(k>d\), on peut, pour tout \(x''\in B\), définir un élément \(\lambda^p(N'',x'')\) par substitution, on encore en utilisant l'unique \(\lambda\)-homomorphisme de \(A\) dans 
